\documentclass{article}

% --- PAQUETES ---
\usepackage[utf8]{inputenc}
\usepackage[spanish]{babel}
\usepackage{geometry}
\usepackage{fancyhdr}
\usepackage{siunitx} % Para la unidad mA

% --- CONFIGURACIÓN DE PÁGINA Y ENCABEZADO ---
\geometry{a4paper, margin=1in, top=1.5in}
\pagestyle{fancy}
\fancyhf{} % Limpiar encabezado y pie de página
\fancyhead[L]{122}
\fancyhead[R]{Prácticas de electricidad}
\renewcommand{\headrulewidth}{0.4pt} % Línea bajo el encabezado

\begin{document}

% Continuación de la lista de preguntas del archivo anterior
\begin{enumerate}
    \setcounter{enumi}{2} % Empezar la numeración en 3
    \item Enunciar la ley de tensión de Kirchhoff. Escribir también la ecuación correspondiente.
    
    \item Explicar los resultados del problema de diseño, apartado 6.
    
    \item ¿Por qué no sería posible diseñar un circuito que diera exactamente una corriente de \SI{5}{mA}?
\end{enumerate}

\vspace{2cm} % Espacio vertical para separar las secciones

\begin{flushleft}
    \textbf{Respuestas a las preguntas de auto-examen}
    
    \vspace{0.5cm}
    
    \begin{enumerate}
        \item 57.
        \item 29.
    \end{enumerate}
\end{flushleft}

\end{document}